\documentclass[10pt,a4paper]{amsart}
\usepackage[margin=19mm]{geometry}
\usepackage{amsmath,amssymb,mathtools}
\usepackage[scr=rsfs,cal=euler]{mathalfa}
\usepackage{xfrac}
\usepackage[unicode,hidelinks,bookmarks=false]{hyperref}
\newcommand{\ie}{i.e.\ }
\newcommand{\cf}{cf.\ }
\newcommand{\resp}{respectively\ }
\newcommand{\Subsection}{\subsection}
\providecommand{\nobreakdash}{\nobreak\hbox{-}}
\setlength{\emergencystretch}{2em}
\multlinegap0pt
\usepackage{amssymb}
\usepackage{mathtools}
\newtheorem{lemma}{Lemma}
\newcommand{\Grd}{\mathrm{Grd}}
\newcommand{\HIT}{\mathrm{HIT}}
\newcommand{\HScert}{\mathrm{HS}^{\mathrm{c}}}
\newcommand{\NZ}{\mathrm{NZ}}
\newcommand{\SigmaB}{\Sigma^B_0}
\newcommand{\Vone}{V^0_1}
\newcommand{\WF}{\mathrm{WF}}
\newcommand{\WFH}{\mathrm{WFH}}
\newcommand{\dWPHPcert}{\mathrm{dWPHP}^{\mathrm{c}}}
\title{Unconditional $\Vone$-independence of a certified hitting-set principle}
\author{Martin Kolář}
\date{Source: arXiv:2608.28813v1 (CC BY 4.0)}
\begin{document}
\maketitle

\noindent\emph{Source and licence.} Unconditional $\Vone$-independence of a certified hitting-set principle. Version arXiv:2608.28813v1 (CC BY 4.0), distributed by arXiv under the Creative Commons Attribution 4.0 International licence, http://creativecommons.org/licenses/by/4.0/, as recorded in the arXiv licence metadata for this exact version and shown on the abstract page https://arxiv.org/abs/2608.28813v1. Full licence notice: \texttt{licenses/arxiv-2608.28813-CC-BY-4.0.txt}.\par\smallskip

\noindent\textit{Editorial scope.} Counted unit: the article's lemma lem:bridge (source0.tex lines 292--294). The article's complete original proof of it is delivered with it (source0.tex lines 296--332, one \texttt{proof} environment of 37 lines). No mathematics is written, repaired or re-derived: the statement and the proof are the article's own lines, in the article's own notation, and no retained line is substituted or re-flowed.  No further source-local context is needed: every cross-reference of every retained line resolves inside the two delivered spans.  Nothing was omitted from any retained span, so the package carries no omission record.  No retained line contains a citation, so the wrapper carries no bibliography and no bibliography entry.  No retained line contains a figure environment, an image-inclusion command or drawing code, so no image is delivered and the package declares no asset dependency.  Editorial formatting only, in addition: an AMS \texttt{amsart} preamble that restates the article's own declarations and macros used by the retained lines, namely the environments \texttt{lemma} on separate counters, together with the standard packages those lines need.  The retained environments print in this document as Lemma 1 in order of appearance, whereas the article prints the same items with the numbering of its own full document.  The complete native source of the article as distributed in the arXiv e-print for this exact version is bundled unchanged as \texttt{sources/208813/source0.tex}.
\begin{lemma}\label{lem:bridge}
$\Vone\vdash\HScert[J]\to\dWPHPcert[J]$.
\end{lemma}

\begin{proof}
Reason in $\Vone$. Fix $N,s,t$ with $\Grd(N,s,t)$ and assume $\HScert[J]$;
let $H$ with $\WFH(H)$ be the hitting tuple it provides. Define
$b:=b(H)\subseteq[N]$ by $b(e(i',j,k)):=\mathrm{bit}(h_{i',j},k)$ for
$(i',j,k)\in[r]\times[n']\times[|q|]$ (a legitimate string by $\SigmaB$-comprehension,
the packing being injective into $[m^3]\subseteq[N]$), and $b(p):=0$ elsewhere.
We claim $b$ witnesses $\dWPHPcert$: no $\WF$ pair certified-computes $b$.

Suppose $\WF(W,Y)$ with $\varepsilon(\cdot;W,Y)=b$ on $[N]$. We derive a
contradiction with the hitting property of $H$.

\emph{(\(\alpha\)) $v_{i'j}(W,Y)=h_{i',j}$ for all $i',j$.} For each bit position
$k<|q|$, $\mathrm{bit}(v_{i'j},k)=\varepsilon(e(i',j,k);W,Y)=b(e(i',j,k))=
\mathrm{bit}(h_{i',j},k)$. Both numbers are $<2^{|q|}$ ($h_{i',j}<q<2^{|q|}$ by
$\WFH$), so binary extensionality on the $|q|$ positions --- a short
($|q|=O(\log N)$) $\SigmaB$-induction --- forces $v_{i'j}=h_{i',j}$.

\emph{(\(\beta\)) nonzeroness.} Take $a:=(2q,\dots,2q)$. Since $q\ge 2^{|q|-1}$
and every $v_{i'j}<2^{|q|}$, we have $a_j>v_{i'j}$ for all $i',j$; hence
$\NZ(a,W,Y)$, and $a\in[2q+1]^{n'}$.

\emph{(\(\gamma\)) $H$ does not hit.} For each $i<r$ instantiate $\HIT$'s outer
index at $i':=i$: by $(\alpha)$, $h_{i,j}=v_{i,j}$ for all $j$, so
$\neg\exists j\,(h_{i,j}\neq v_{i,j})$, i.e.\ $\neg\HIT(i,H,W,Y)$. As this holds
for every $i<r$, we have $\neg\exists i<r\,\HIT(i,H,W,Y)$.

Now $\WF(W,Y)\wedge(\exists a\in[2q+1]^{n'}\,\NZ(a,W,Y))$ holds by $(\beta)$,
so the inner implication of $\HScert$ yields $\exists i<r\,\HIT(i,H,W,Y)$,
contradicting $(\gamma)$. Hence no $\WF$ pair certified-computes $b$, and
$\dWPHPcert[J]$ holds.

Each step uses only $\SigmaB$-comprehension, short $\SigmaB$-induction on
$|q|=O(\log N)$ bit positions, and $\Delta_0$ arithmetic; the algebraic
evaluation of $A_{(W,Y)}$ is never performed --- the unfolded forms $\NZ,\HIT$
replace it --- so the $\mathrm{TC}^0$-hardness of that evaluation is avoided.
No amplification and no counting enters.
\end{proof}

\end{document}
